\chapter{Aplicaciones Django}
\label{chap:aplicaciones-django}

\section{Organizacion del dominio por aplicaciones}
\label{sec:apps-organizacion}

La plataforma divide su dominio en cinco aplicaciones locales: \archivo{apps.common}, \archivo{apps.core}, \archivo{apps.participants}, \archivo{apps.tournament} y \archivo{apps.invitations}. Esta separacion no responde solo a una convencion de carpetas; define como fluye la informacion desde la entrada publica hasta la competencia y las comunicaciones.

El sitio publico presenta informacion de la edicion activa. El registro captura datos de equipos y participantes. La confirmacion de asistencia transforma registros en equipos operativos. El modulo de torneo usa esos equipos para construir competencias por division. El modulo de comunicaciones usa registros, destinatarios, plantillas y logs para controlar envios y solicitudes. \archivo{apps.common} aporta una base reutilizable para registrar fechas de creacion y actualizacion en modelos de dominio.

\begin{figure}[H]
  \centering
  \fbox{\parbox{0.88\textwidth}{\centering Marcador de diagrama: dependencias funcionales entre aplicaciones Django. Fuente prevista: DIA-03 en \archivo{planificacion/06_inventario_diagramas.md}.}}
  \caption{Marcador para diagrama de dependencias entre aplicaciones.}
  \label{fig:apps-dependencias}
\end{figure}

\section{Mapa de responsabilidades}
\label{sec:apps-mapa-responsabilidades}

\begin{longtable}{p{3.2cm}p{3.6cm}p{6.2cm}}
\toprule
Aplicacion & Configuracion & Responsabilidad principal \\
\midrule
\archivo{apps.common} & \clase{CommonConfig} & Modelo abstracto de timestamps para entidades persistentes. \\
\archivo{apps.core} & \clase{CoreConfig} & Sitio publico: home, reglas publicadas y patrocinadores. \\
\archivo{apps.participants} & \clase{ParticipantsConfig} & Registros, participantes, asistencia, instituciones, equipos y sincronizacion operativa. \\
\archivo{apps.tournament} & \clase{TournamentConfig} & Ediciones, reglas, competencias por division, grupos, batallas, estados, recomendaciones y panel interno. \\
\archivo{apps.invitations} & \clase{InvitationsConfig} & Plantillas DOCX, destinatarios, lotes, logs, solicitudes de asistencia y envios de comunicacion. \\
\bottomrule
\end{longtable}

La tabla resume responsabilidades, pero el punto critico esta en las dependencias. \archivo{apps.participants} depende de \archivo{apps.tournament} para asociar registros a una edicion activa y para inicializar competencias cuando se confirma asistencia. \archivo{apps.tournament} depende de \archivo{apps.participants} porque las competencias operan sobre \clase{Team} y \clase{TeamMember}. \archivo{apps.invitations} consulta registros de participantes para comunicaciones y solicitudes de asistencia. \archivo{apps.core} consume informacion de torneo y participantes para presentar estado publico del evento.

\section{apps.common}
\label{sec:apps-common}

\archivo{apps.common} contiene \clase{TimeStampedModel}, clase abstracta con \variable{created_at} y \variable{updated_at}. La aplicacion no implementa flujo funcional propio, pero reduce duplicacion en modelos de participantes, torneo e invitaciones.

Su impacto de mantenimiento es bajo pero transversal. Si se cambia \clase{TimeStampedModel}, el cambio afecta entidades de varios dominios. Por eso esta aplicacion debe permanecer pequeña y estable. Convertirla en un contenedor de utilidades sin relacion clara le haria perder su valor como base comun.

\section{apps.core}
\label{sec:apps-core}

\archivo{apps.core} implementa la entrada publica de la plataforma. Sus vistas consultan la edicion activa, reglas publicadas, conteos de registros y catalogo de patrocinadores. No registra participantes ni modifica competencias; su funcion es presentar informacion verificable y enlazar al flujo publico.

\begin{longtable}{p{3.4cm}p{3.8cm}p{5.8cm}}
\toprule
Modulo & Elementos principales & Funcion dentro del flujo \\
\midrule
\archivo{views.py} & \funcion{sponsor_catalog}, \funcion{home}, \funcion{rules_page}, \funcion{sponsors_page} & Construye contexto publico desde modelos de torneo y participantes. \\
\archivo{urls.py} & \ruta{/}, \ruta{/reglas/}, \ruta{/patrocinadores/} & Expone rutas publicas sin autenticacion. \\
\archivo{templates/core/} & \archivo{base.html}, \archivo{home.html}, \archivo{rules.html}, \archivo{sponsors.html} & Renderiza informacion publica y navegacion hacia registro. \\
\bottomrule
\end{longtable}

La dependencia mas relevante de \archivo{apps.core} es de lectura. \funcion{home} consulta \clase{TournamentEdition}, \clase{RuleSection}, \clase{SchoolRegistration} y \clase{UniversityRegistration}. Esta lectura permite mostrar estado de la edicion sin duplicar datos en otra tabla. Si se cambian nombres de estados o reglas de edicion activa, el sitio publico puede cambiar de contenido aunque no se modifique \archivo{apps.core}.

\section{apps.participants}
\label{sec:apps-participants}

\archivo{apps.participants} captura y transforma informacion de participantes. Su flujo empieza con un formulario publico y termina, cuando hay confirmacion de asistencia, en entidades operativas para torneo.

\begin{longtable}{p{3.4cm}p{4cm}p{5.6cm}}
\toprule
Modulo & Responsabilidad & Interaccion \\
\midrule
\archivo{models.py} & Registros, participantes, instituciones, equipos y miembros & Define datos que conectan inscripcion publica y competencia. \\
\archivo{forms.py} & Validaciones de registro escolar y universitario & Bloquea duplicados de robot, correo y documentos en la edicion activa. \\
\archivo{views.py} & Registro, exito y confirmacion por token & Orquesta formulario, servicios de correo y confirmacion de asistencia. \\
\archivo{services.py} & Correos, confirmacion y sincronizacion a equipo & Actualiza registros y crea/actualiza \clase{Institution}, \clase{Team} y \clase{TeamMember}. \\
\archivo{admin.py} & Administracion de registros y equipos & Permite acciones de solicitud de asistencia y sincronizacion de confirmados. \\
\archivo{templates/participants/} & Formularios y pantallas de confirmacion & Presenta flujo publico al usuario externo. \\
\bottomrule
\end{longtable}

El flujo de informacion funciona asi: una vista de registro instancia \clase{SchoolRegistrationForm} o \clase{UniversityRegistrationForm}; el formulario valida la edicion activa y duplicados; al guardar crea un registro y sus participantes; la vista llama \funcion{send_registration_received_email}; despues, la asistencia se confirma por token en \funcion{attendance_confirmation}; el servicio \funcion{confirm_attendance} aprueba el registro, sincroniza equipo y puede inicializar competencia si corresponde.

Esta app es el puente entre usuarios publicos y competencia. Cambiar sus validaciones impacta el conjunto de equipos que llega a \archivo{apps.tournament}. Cambiar \funcion{sync_registration_to_team} impacta relaciones de \clase{Team}, \clase{TeamMember}, estados de competencia y panel interno.

\section{apps.tournament}
\label{sec:apps-tournament}

\archivo{apps.tournament} concentra el nucleo competitivo. No solo almacena fases; tambien calcula perfiles, construye grupos, registra resultados, crea repechajes, permite fases manuales, sincroniza estados y cierra podios.

\begin{longtable}{p{3.6cm}p{4.2cm}p{5.2cm}}
\toprule
Modulo & Responsabilidad & Impacto \\
\midrule
\archivo{models.py} & Ediciones, reglas, competencias, grupos, batallas, estados, historial y matches genericos & Define el estado persistente del torneo. \\
\archivo{formats.py} & Perfiles automaticos segun cantidad de equipos & Determina estructura inicial de competencia. \\
\archivo{planner.py} & Planificacion y distribucion de slots & Apoya distribucion operativa. \\
\archivo{recommendations.py} & Recomendaciones de fase, repechaje y tercer lugar & Alimenta decisiones del panel interno. \\
\archivo{services.py} & Regla de negocio principal de torneo & Modifica grupos, batallas, estados, historial y podio. \\
\archivo{views.py} & Vista publica y panel interno & Convierte acciones HTTP en llamadas a servicios. \\
\archivo{urls.py} & Rutas de torneo y control & Expone entrada publica y operaciones internas protegidas. \\
\archivo{admin.py} & Administracion de modelos de torneo & Permite inspeccion y ajustes administrativos. \\
\archivo{tests.py} & Pruebas de distribuciones manuales & Cubre parte critica de fases manuales. \\
\bottomrule
\end{longtable}

La dependencia con \archivo{apps.participants} es estructural: \clase{DivisionCompetition} se construye por edicion y division, pero los participantes reales de competencia son \clase{Team}. Los grupos y batallas referencian equipos; los estados de competencia siguen el avance de cada equipo; el historial registra cambios por equipo. Por tanto, el torneo no puede entenderse sin los modelos operativos de participantes.

El panel interno no escribe directamente toda la logica. Sus vistas llaman servicios como \funcion{initialize_competition}, \funcion{set_group_qualifiers}, \funcion{set_battle_winner}, \funcion{set_battle_qualifiers}, \funcion{materialize_repechage_stage}, \funcion{generate_manual_phase_proposal}, \funcion{confirm_manual_phase_proposal} y \funcion{save_final_podium}. Esta organizacion reduce duplicacion entre vistas, pero concentra riesgo en \archivo{services.py}.

\section{apps.invitations}
\label{sec:apps-invitations}

\archivo{apps.invitations} gestiona comunicaciones y solicitudes. Su responsabilidad es transformar plantillas y destinatarios en lotes trazables, con controles para simulacion, prueba controlada y envio oficial.

\begin{longtable}{p{3.4cm}p{4.2cm}p{5.4cm}}
\toprule
Modulo & Responsabilidad & Interaccion \\
\midrule
\archivo{models.py} & Plantillas, destinatarios, lotes y logs & Persiste la trazabilidad de comunicaciones. \\
\archivo{forms.py} & Validacion de plantillas, destinatarios y lotes & Bloquea archivos invalidos, plantillas inactivas y envios reales sin confirmacion. \\
\archivo{services.py} & DOCX, PDF, XLSX, correo y logs & Ejecuta transformaciones y registra resultados por destinatario. \\
\archivo{views.py} & Dashboard, plantillas, invitaciones, confirmaciones e historial & Orquesta formularios y servicios desde rutas internas. \\
\archivo{admin.py} & Administracion de comunicaciones & Permite inspeccion de plantillas, lotes, destinatarios y logs. \\
\archivo{tests.py} & Pruebas de PDF y correo & Valida componentes sensibles del modulo. \\
\bottomrule
\end{longtable}

La app interactua con \archivo{apps.participants} porque las solicitudes de asistencia se basan en registros escolares y universitarios. Tambien interactua con configuracion global: SMTP determina si puede enviar correos reales, y \variable{LIBREOFFICE_BINARY} o PATH determinan si puede convertir documentos a PDF. Por eso sus fallos no siempre nacen del codigo de la app; pueden surgir de configuracion de entorno.

\section{Flujo de informacion entre aplicaciones}
\label{sec:apps-flujo-informacion}

El flujo principal del dominio es secuencial, aunque cada aplicacion mantenga su responsabilidad:

\begin{enumerate}
  \item \archivo{apps.core} presenta informacion de la edicion activa y dirige al registro.
  \item \archivo{apps.participants} crea el registro y participantes asociados.
  \item \archivo{apps.invitations} puede generar comunicaciones o solicitudes de asistencia sobre esos registros.
  \item \archivo{apps.participants} confirma asistencia por token y sincroniza a equipos.
  \item \archivo{apps.tournament} consume equipos para construir y operar competencias por division.
  \item \archivo{apps.tournament} registra resultados, estados e historial.
  \item \archivo{apps.core} y vistas internas reflejan datos actualizados para consulta publica u operativa.
\end{enumerate}

\begin{figure}[H]
  \centering
  \fbox{\parbox{0.88\textwidth}{\centering Marcador de diagrama: flujo de informacion entre aplicaciones desde registro hasta competencia. Fuente prevista: DIA-07 y DIA-08 en \archivo{planificacion/06_inventario_diagramas.md}.}}
  \caption{Marcador para flujo de informacion entre aplicaciones.}
  \label{fig:apps-flujo-informacion}
\end{figure}

\section{Django Admin como superficie operativa}
\label{sec:apps-admin}

Django Admin forma parte de la arquitectura operativa. \archivo{apps.participants/admin.py} incluye acciones que pueden enviar solicitudes de asistencia o sincronizar registros confirmados. \archivo{apps.tournament/admin.py} permite inspeccionar entidades del torneo. \archivo{apps.invitations/admin.py} permite revisar plantillas, destinatarios, lotes y logs.

El valor de Admin es la capacidad de inspeccion y correccion. El riesgo es que algunas acciones tienen efectos reales, especialmente correos y sincronizacion de equipos. Por esa razon, las acciones administrativas deben documentarse y probarse como parte de mantenimiento, no tratarse como herramientas auxiliares sin impacto.

\section{Consideraciones de mantenimiento}
\label{sec:apps-mantenimiento}

\begin{longtable}{p{4cm}p{9cm}}
\toprule
Aspecto & Consideracion \\
\midrule
Archivos relacionados & \archivo{apps/*/models.py}, \archivo{forms.py}, \archivo{views.py}, \archivo{services.py}, \archivo{urls.py}, \archivo{admin.py} y templates de cada app. \\
Riesgo al modificar & Cambios en participantes pueden alterar equipos; cambios en torneo pueden afectar resultados e historial; cambios en invitaciones pueden enviar comunicaciones incorrectas. \\
Dependencias & \archivo{participants} y \archivo{tournament} estan acopladas por \clase{Team}; \archivo{invitations} depende de registros y configuracion SMTP/PDF; \archivo{core} consume datos de lectura. \\
Pruebas recomendadas & Registro escolar y universitario, confirmacion por token, sincronizacion a equipo, generacion de competencia, acciones de panel, envio simulado y prueba controlada. \\
Impacto sobre otros modulos & Modificar estados, nombres de division o estructura de equipos impacta vistas, servicios, templates, Admin y reportes de comunicacion. \\
Buenas practicas & Mantener regla de negocio en servicios, evitar duplicar validaciones entre vistas, documentar nuevas acciones Admin y actualizar pruebas al cambiar flujos. \\
\bottomrule
\end{longtable}
